\begin{definition}[The published special-unitary Tamagawa input]
Let $F$ and $E$ be number fields, with a specified $F$-algebra structure
on $E$ and $\sigma\in\operatorname{Aut}_F(E)$ satisfying the imported
CM-extension condition $\operatorname{IsCMExt}(F,E,\sigma)$.
Let $M\ge2$, let $J\in M_M(E)$ satisfy
$\operatorname{IsHermND}(\sigma,J)$, and let
$\Lambda\subset E^M$ be a finitely generated $\mathcal O_E$-submodule
with $E\Lambda=E^M$. Put $d=M^2-1$ and
$G(\mathbb A_F)=\operatorname{SUAdelic}(\sigma,J)$, with its specified
rational subgroup $\Gamma=\operatorname{SURat}(\sigma,J)$.

The local scalars used in this input are the ones in
\texttt{GlobalTheory.TamSU.locVol}, with the following precise
normalization. Let
\[
 K_\Lambda=\{h\in G(\mathbb A_F):
                 h_{\mathrm{fin}}\widehat\Lambda=\widehat\Lambda\}.
\]
For a nonzero prime ideal $v\subset\mathcal O_F$ and $k\ge0$ let
$c_v(k)$ be the source's natural-valued cardinality (\texttt{Nat.card},
which is zero for an infinite quotient) of the quotient of $K_\Lambda$ by
\[
 h\sim_{v,k}h'
 \quad\Longleftrightarrow\quad
 (h_{\mathrm{fin}}-h'_{\mathrm{fin}})\widehat\Lambda
       \subset\widehat{((v^k\mathcal O_E)\Lambda)}.
\]
Set $q_v=N_{F/\mathbb Q}(v)$ and
\[
 a_v=\operatorname{limUnder}_{k\to\infty}
                         \frac{c_v(k)}{q_v^{kd}} .
\]
Here $\operatorname{limUnder}$ is the exact limit-selection operation
used by the imported definition; convergence of each sequence is not
an extra assertion of this module.

Define the integral Lie lattice
\[
 \mathfrak g_\Lambda=
 \{X\in M_M(E):\operatorname{tr}X=0,\ 
        (\sigma X)^tJ+JX=0,\ X\Lambda\subset\Lambda\}.
\]
Choose a finite CM type $\Psi\subset\operatorname{Hom}(E,\mathbb C)$:
for every $\tau$, $\tau\in\Psi$ if and only if
$\overline\tau\notin\Psi$.
Choose maps
$\theta_\rho:\mathbb A_{F,\infty}\otimes_F E\to\mathbb C$
for $\rho\in\Psi$ with the imported compatibility
$\operatorname{IsInfEmb}(\rho,\theta_\rho)$.
Choose $e_1,\ldots,e_d\in\mathfrak g_\Lambda$ linearly independent
over $F$. Write $L_e=\sum_i\mathcal O_F e_i$ and
$I_e=\operatorname{relIndex}(L_e,\mathfrak g_\Lambda)$,
in the source's order of the two subgroups.

Let $\nu_{\Psi,e}$ be the pushforward of Lebesgue measure on
$\mathbb R^{\Psi\times\{1,\ldots,d\}}$ by
\[
 (c_{\rho i})\longmapsto
       \left(\sum_{i=1}^d c_{\rho i}\rho(e_{i,ab})\right)_{\rho,a,b}.
\]
This is exactly \texttt{SUVol.formLie}; put
$B_\epsilon=\{X:\lvert X_{\rho,ab}\rvert<\epsilon
\text{ for every }\rho,a,b\}$.
For a measure $\mu$ on $G(\mathbb A_F)$,
$\operatorname{IsTamagawa}(\sigma,J,\Lambda,\Psi,\theta,e,\mu)$
means that $\mu$ is Haar and, as $\epsilon\downarrow0$,
\[
 \frac{\mu\{h\in K_\Lambda:
       |h_{\rho,ab}-\delta_{ab}|<\epsilon
                  \text{ for every }\rho,a,b\}}
      {\nu_{\Psi,e}(B_\epsilon)}
 \longrightarrow
 \operatorname{ofReal}\left(
     |\operatorname{disc}F|^{-d/2} I_e\prod_v' a_v\right).
\]
Here $h_\rho$ is the actual archimedean matrix obtained using
$\theta_\rho$, and the quotient and limit are in $[0,\infty]$,
as in \texttt{IsTamagawa}.

The proposition \texttt{Published.TamagawaNumberSU} asserts,
for every $F,E,\sigma,M,J,\Lambda$ above, both:
\begin{enumerate}
\item The real family $(a_v-1)_v$, indexed by all nonzero prime
ideals of $\mathcal O_F$, is summable.
\item For every $\Psi,\theta,e$ as above, every measure $\mu$
satisfying this exact $\operatorname{IsTamagawa}$ normalization,
and every fundamental domain $\mathcal F$ for the specified
rational subgroup $\Gamma\subset G(\mathbb A_F)$,
\[
                              \mu(\mathcal F)=1.
\]
\end{enumerate}
\end{definition}

\begin{remark}
This file defines a proposition used as a published-theorem input;
it contains no proof of that proposition. In particular it does not
assert existence of a normalized measure or a fundamental domain.
The finite-place summability clause, the rank restriction $M\ge2$,
the full-lattice hypothesis, and the displayed Lie-lattice and
discriminant normalization are part of the input, not optional
conventions. The auxiliary formulas above were read from
\texttt{Definitions.Divisors.IsTamagawa} and
\texttt{PublishedTheorems.Divisors$_2$.EmbGLMemRealSU}.
\end{remark}

